% =============================================
% File: appendix.tex
% Phụ lục
% =============================================

\appendix
\renewcommand{\chaptername}{Phụ lục}

% ============================================================
% Phụ lục A — Mã nguồn bộ điều khiển PI trên ESP32
% ============================================================
\chapter{MÃ NGUỒN BỘ ĐIỀU KHIỂN PI TRÊN ESP32}
\label{app:esp32_code}

Đoạn mã sau trình bày cấu trúc chính của chương trình điều khiển PI vận tốc
triển khai trên vi điều khiển ESP32 WROOM-32 (Arduino framework).

\begin{verbatim}
// --- Cấu hình phần cứng ---
#define ENC_A 18
#define ENC_B 19
#define PWM_PIN 25
#define DIR_PIN1 26
#define DIR_PIN2 27

volatile long encoderCount = 0;

// Ngắt đọc encoder
void IRAM_ATTR encoderISR() {
    encoderCount += (digitalRead(ENC_B) == HIGH) ? 1 : -1;
}

// Tính vận tốc (RPM)
float VelCalculate(long &posNow, float dt) {
    float vel = (posNow / 200.0f) / dt * 60.0f;
    posNow = 0;
    return vel;
}

// Biến PI
float Kp = 0.185436f;
float Ki = 5.288458f;
float integral = 0.0f;
const float DT = 0.01f;
const float INT_LIMIT = 1000.0f;

// Vòng điều khiển chính (gọi mỗi 10 ms)
void piControl(float velReq, float velCur) {
    float err = velReq - velCur;
    integral += err * DT;
    integral = constrain(integral, -INT_LIMIT, INT_LIMIT);
    float u = Kp * err + Ki * integral;
    u = constrain(u, -255, 255);
    // Ghi PWM ra động cơ
    int pwmVal = (int)abs(u);
    digitalWrite(DIR_PIN1, u >= 0 ? HIGH : LOW);
    digitalWrite(DIR_PIN2, u >= 0 ? LOW  : HIGH);
    ledcWrite(0, pwmVal);
}
\end{verbatim}


% ============================================================
% Phụ lục B — Mã nguồn huấn luyện mạng ANN (Python)
% ============================================================
\chapter{MÃ NGUỒN HUẤN LUYỆN MẠNG ANN}
\label{app:ann_code}

Đoạn mã Python dưới đây trình bày kiến trúc lớp và vòng huấn luyện của
mạng ANN được xây dựng từ đầu (không dùng TensorFlow/PyTorch).

\begin{verbatim}
import numpy as np

# ---- Lớp Dense với L2 regularization ----
class DenseL2:
    def __init__(self, n_in, n_out, lam=5e-4):
        self.W = np.random.randn(n_in, n_out) * 0.01
        self.b = np.zeros((1, n_out))
        self.lam = lam

    def forward(self, X):
        self.X = X
        return X @ self.W + self.b

    def backward(self, dOut, lr):
        dW = self.X.T @ dOut + self.lam * self.W
        db = dOut.sum(axis=0, keepdims=True)
        dX = dOut @ self.W.T
        self.W -= lr * dW
        self.b -= lr * db
        return dX

# ---- Kiến trúc: 2 -> 8 -> 16 -> 2 ----
# Input: [error_vel, delta_error_vel]
# Output: [Kp, Ki]
#
# layers = [DenseL2(2,8), ReLU, Dropout(0.1),
#           DenseL2(8,16), ReLU, Dropout(0.1),
#           DenseL2(16,2), Linear]
#
# Optimizer: Adam  lr=2e-3  decay=3e-4
# Loss: MSE   Epochs: 1000
\end{verbatim}


% ============================================================
% Phụ lục C — Thông số tối ưu PI tại 9 mức setpoint
% ============================================================
\chapter{BẢNG THÔNG SỐ PI TỐI ƯU}
\label{app:pi_table}

Bảng~\ref{tab:pi_optimal} tổng hợp bộ tham số $K_p^*$, $K_i^*$ tối ưu
được xác định bằng giải thuật di truyền (GA) tại 9 mức setpoint.

\begin{table}[H]
\centering
\caption{Bộ tham số PI tối ưu tại các mức setpoint}
\label{tab:pi_optimal}
\renewcommand{\arraystretch}{1.3}
\begin{tabular}{|c|c|c|}
\hline
\textbf{Setpoint (RPM)} & \textbf{$K_p^*$} & \textbf{$K_i^*$} \\
\hline
250  & 0,185436 & 5,288458 \\
500  & --       & --       \\
750  & --       & --       \\
1000 & --       & --       \\
1250 & --       & --       \\
1500 & --       & --       \\
1750 & --       & --       \\
2000 & --       & --       \\
2250 & --       & --       \\
\hline
\end{tabular}

\vspace{0.3cm}
\textit{Ghi chú: Các ô ``--'' cần bổ sung từ kết quả thực nghiệm.}
\end{table}


% ============================================================
% Phụ lục D — Đặc tính động học cơ cấu 5-bar Parallel
% ============================================================
\chapter{ĐỘNG HỌC CƠ CẤU 5-BAR PARALLEL}
\label{app:fivebar}

\section*{Thông số hình học}
\begin{itemize}
    \item Khoảng cách hai tâm khớp chủ động: $A_1A_2 = 105$ mm
    \item Chiều dài khâu chủ động: $L_{act} = 100$ mm
    \item Chiều dài khâu bị động: $L_{pas} = 180$ mm
\end{itemize}

\section*{Động học thuận (FK)}
Cho trước hai góc khớp $\theta_1$, $\theta_2$, tọa độ điểm cuối $P(x_P, y_P)$
được xác định qua hệ phương trình:
\begin{align}
    B_1 &= (L_{act}\cos\theta_1,\; L_{act}\sin\theta_1) \\
    B_2 &= (A_1A_2 + L_{act}\cos\theta_2,\; L_{act}\sin\theta_2)
\end{align}
$P$ là giao điểm của hai đường tròn bán kính $L_{pas}$ tâm $B_1$ và $B_2$.

\section*{Động học nghịch (IK)}
Cho trước tọa độ $P(x_P, y_P)$, tính ngược ra $\theta_1$, $\theta_2$:
\begin{align}
    \theta_1 &= \text{atan2}(B_{1y} - 0,\; B_{1x} - 0) \\
    \theta_2 &= \text{atan2}(B_{2y} - 0,\; B_{2x} - A_1A_2)
\end{align}

\section*{Ma trận Jacobian}
Ma trận Jacobian $J(\theta)$ liên hệ vận tốc khớp với vận tốc điểm cuối:
\begin{equation}
    \begin{pmatrix} \dot{x}_P \\ \dot{y}_P \end{pmatrix}
    = J(\theta)
    \begin{pmatrix} \dot{\theta}_1 \\ \dot{\theta}_2 \end{pmatrix}
\end{equation}
